\section{Introduction}\label{sec:introduction}

Many existence problems in arithmetic begin with a comparison between global
objects and their local avatars. The central question is whether compatible
local data can be assembled into a global solution. In the setting considered
here, the obstruction is encoded by a finite set $\Omega$ and a map
\begin{equation}\label{eq:obstruction-map}
  \operatorname{ob}\colon \mathcal{X}\longrightarrow \prod_{v\in\Omega}
  \mathcal{X}_v.
\end{equation}
Our goal is to isolate hypotheses under which the image of~\eqref{eq:obstruction-map}
is large enough to detect global solvability.

The main result, Theorem~\ref{thm:main}, has two useful features. First, its
hypotheses are entirely local away from the finite set $\Omega$. Second, the
argument gives a lower bound rather than existence alone. This separation is
useful when the same geometric input is paired with different counting
methods.

\subsection{Outline of the paper}

Section~\ref{sec:framework} fixes notation and records the local hypotheses.
Section~\ref{sec:main-result} states the main theorem and organizes its proof
into two lemmas. Section~\ref{sec:questions} collects variants and open
questions.
